\documentclass[../main.tex]{subfiles}

\begin{document}

\chapter{Conclusion}\label{ch:conclusion}

This thesis presented an application layer monitoring tool for the Shibboleth IdP that detects 
documented vulnerabilities, configuration weaknesses and anomalous authentication behaviour by 
combining three separate sources of information from within the identity provider. 
The collected data is transmitted to an external backend, where it is stored, correlated and 
evaluated against a catalogue of thirteen detection rules.

The three sources proved complementary, and the reason is more specific than their number. 
The audit log records the outcome of a completed transaction, the configuration states what was 
intended for the partner involved, and the interceptor observes the incoming message while it is 
being processed. 
What each source can see follows from its position in the processing sequence, which is why their 
combination permits comparisons that no single source supports on its own. 
The comparison between a declared configuration and an observed event is the clearest example, 
since a login without multi factor authentication is unremarkable until it is related to a service 
provider that requires it.

The detection rules are derived rather than assembled from experience. 
Eleven of the thirteen rules can be traced to vulnerabilities and attack classes documented in the 
literature or to normative requirements for SAML deployments, and the two that cannot are named as 
such. 
All thirteen are verified against a real Elasticsearch instance, and five additionally through the 
complete pipeline from a browser login to the resulting alert.

The interceptor introduced no measurable overhead in the asynchronous transmission mode, with the 
work performed inside it amounting to approximately 0,3\% of the login duration. 
The approach can therefore be integrated into the authentication workflow of a central 
authentication service without a perceptible effect on the logins it observes.

A recurring observation across the evaluation deserves mention, since it bears on the premise of 
the work. 
In each of the three evaluations, the measurement uncovered a defect that operation itself had not 
revealed. 
The system continued to run in every case, alerts continued to be produced and the dashboard 
continued to be populated. 
The defects became apparent only when the behaviour was compared against an expectation defined in 
advance. 
That a system appears to function is therefore not evidence that it does, which is the same 
argument that motivates observing an identity provider from within in the first place.

The study also identified limitations that follow from the observation point rather than from the 
effort invested. 
Several federation level attack scenarios occur at a different system layer, involve a different 
party, or are observable only through their preconditions. 
One case is more fundamental, since a malicious identity provider cannot be detected by a system 
running inside it, where observer and adversary coincide. 
The severity classification makes a further boundary explicit, as none of the thirteen rules 
reaches the highest risk level, which presupposes a compromise of the signing key. 
The Sunburst incident that motivated this work would consequently not have been detected by the 
system presented here, although the configuration weaknesses that make such an attack easier to 
carry out and harder to notice are within its reach.

What this thesis contributes is therefore not a new architecture for collecting and correlating 
security events, since that is established practice. 
It is the demonstration that an identity provider holds security relevant information which no 
log consuming system can obtain, an analysis of which properties are observable from which 
position, and an implementation that makes this information available at a cost that can be 
stated in numbers rather than assumed.

\end{document}
